\chapter{Проектирование архитектуры и подбор инструментов системы}\label{ch:2}

% Как вообще пришли к разработке своего микросервиса?
% Основные функциональные возможности разрабатываемой системы
% Сущности и что они означают (ER диаграммы)
% Нужно хранилище ключей для RSA. Выбираем Vault, рассказать что такое Vault. Причина – уже используется в компании, есть инфраструктура, лучшее решение на рынке для хранения секретов.
% В системе гаранитровано будем хранить данные (см сущности). Нужен persistence, потеря информациии недопустима. Выбираем PostgreSQL – есть экспертиза, готова инфраструктура, подходит для транзакционной нагрузки, запросы к базе – простые селекты, не больше.
% Выдача токенов должна быть централизована, ротация должна происходить автоматически. Делаем свой микросервис. Стандарт в компании – C#. Перечислить преимущества языка, развитую экосистему библиотек от Microsoft, особенно по работе с авторизацией.
% Архитектура типового микросервиса в компании (Model, Services, Worker, Migrator). Описать типичную архитектуру с репозиториями, сервисами, контроллерами.
% Обязательные интеграционные тесты. Testcontainers с воссозданием всей инфраструктуры в контейнерах.
% Мониторинг через библиотеку prometheus-net, формат Prometheus.
% Хотим хранить клиентов авторизации кодом. Выбираем написать собственные Terraform Provider (Go) и Terraform Module (HCL). Объяснить почему такой выбор.
% Практически все микросервисы используют Mindbox.Authentication для работы с авторизацией. Забираем библиотеку себе на поддержку, существенно переделываем для новой схемы работы авторизации
% Хотим сделать микросервис надежным и быстрым – добавляем распределенный кэш (Dragonfly). Рассказать почему не Redis, преимущества, кратко. Показать схему получения данных L1 (in memory) -> L2 (distributed cache) -> L3 (база данных)
% Не хотим стать единой точкой отказа – добавляем асинхронный механизм доставки данных (JWKs, токены) до микросервсов. Рассмотреть реализации через агент и через оператор. Показать почему оператор лучше, пишем на Go и Operator SDK, тоже пояснить почему эти технологии.
% Во всех пунктах надо опираться на технические требования к системе из первой главы.
% В выводе надо все резюмировать, показать какие требования закрываем точно, какие гипотетически, какие не закрываем. Которые не закрываем объяснить почему готовы делать без них.
